\section{Objetivos de la investigación}

\subsection{Objetivo general}

Implementar un modelo crítico, \textquotedbl{}Critic-EmpRAG1.0\textquotedbl{}, sobre Arquitecturas de autorreflexión y RAG aplicadas a modelos de lenguaje pequeños para la mejora de la productividad y el acceso al Sistema de soporte de conocimiento empresarial interno de las áreas de TI y desarrollo en Dualbiz SRL.

\subsection{Objetivos específicos}

\begin{enumerate}
	
	\item Sistematizar la literatura acerca de los métodos, técnicas y despliegues de Arquitecturas de autorreflexión y generación aumentada por recuperación (RAG); con énfasis en su aplicabilidad a modelos de lenguaje pequeños, contextualizando a tareas sobre bases de conocimiento empresarial.
	
	\item Diagnosticar el estado actual de los procesos en el acceso al Sistema de soporte de conocimiento empresarial interno en las áreas de TI y desarrollo en Dualbiz SRL; para la indentificación de los problemas actuales suscitados, sus restricciones de privacidad de datos y de infraestructura de cómputo.
	
	\item Diseñar el Modelo crítico, Critic-EmpRAG1.0, y construir su prototipo funcional sobre modelos de lenguaje pequeños, para la mejora de la productividad y el acceso al Sistema de soporte de conocimiento empresarial interno de las áreas de TI y desarrollo en Dualbiz SRL.
	
	\item Validar el Modelo crítico, Critic-EmpRAG1.0, teóricamente mediante la consulta a expertos y empíricamente en dos niveles ---organizacional, sobre la productividad del proceso de consulta, y computacional, sobre el acceso al conocimiento frente a configuraciones alternativas del sistema---, en la mejora de la productividad y el acceso al Sistema de soporte de conocimiento empresarial interno de las áreas de TI y desarrollo en Dualbiz SRL, con restricciones de infraestructura de cómputo y de acceso por privacidad de la información.	 
		
\end{enumerate}

